\documentclass[11pt]{article}
\usepackage[review]{acl}

\usepackage{times}
\usepackage{latexsym}
\usepackage[T1]{fontenc}
\usepackage[utf8]{inputenc}
\usepackage{microtype}
\usepackage{inconsolata}
\usepackage{graphicx}
\usepackage{booktabs}
\usepackage{multirow}
\usepackage{array}
\usepackage{amsmath}
\usepackage{tabularx}

\newcommand{\PLACEHOLDER}[1]{\textbf{[#1]}}

\title{SafeSteer-IN: Inference-Time Safety Steering for Indic LLMs with\
Contrastive Activation Addition}

\author{Anonymous ACL Submission}

\begin{document}
\maketitle

\begin{abstract}
Large language models are being rapidly deployed for Indian users, but safety tooling remains largely English-centric and often misses culturally specific harms such as caste discrimination, communal incitement, and code-mixed abuse. We present \textbf{SafeSteer-IN}, a plug-in inference-time safety layer for Indic LLMs based on Contrastive Activation Addition (CAA). Our work contributes two components: (i) an India-grounded safety taxonomy spanning 8 harm categories across 6 Indic language settings, and (ii) an end-to-end steering pipeline that routes prompts by language/category and applies slice-specific activation steering without model weight updates. The system integrates dataset construction, steering-vector extraction, alpha calibration, and evaluation in a reproducible workflow. Preliminary experiments on Sarvam-1 and OpenHathi indicate promising safety steering behavior in selected Hindi/Tamil slices while preserving generation quality, with quantitative values left as placeholders pending final runs. We position this as an early but concrete step toward culturally grounded safety mechanisms for multilingual LLM deployment.
\end{abstract}

\section{Introduction}
India now has more than 600M internet users across multiple scripts and language families, while Indic LLM deployment is accelerating through systems such as Sarvam-1, OpenHathi, Krutrim, and newer large Sarvam releases. Yet most mainstream safety pipelines are still built around English-centric moderation data and taxonomies. This creates a practical mismatch: harmful requests in Indian sociolinguistic contexts are often subtle, code-mixed, and culturally specific, but are evaluated by filters that were not designed for these phenomena.

A representative failure mode is a Hindi chatbot responding to a softly phrased communal incitement query. Because the prompt may not contain canonical English hate markers, an English-trained safety layer can under-trigger and allow inflammatory content generation. Similar failures appear for caste-targeted slurs and mixed-script toxicity (e.g., Hinglish), where lexical and cultural cues differ from common Western benchmark assumptions.

In this work, we argue that culturally grounded safety for Indic LLMs requires both \emph{taxonomy-level adaptation} and \emph{mechanism-level intervention}. We therefore build SafeSteer-IN around Contrastive Activation Addition (CAA), a lightweight inference-time approach that does not fine-tune model weights.

Our contributions are:
\begin{itemize}
    \item \textbf{India-specific safety taxonomy:} 8 harm categories over 6 Indic language settings, including culturally grounded categories underrepresented in existing multilingual safety benchmarks.
    \item \textbf{First CAA application to Indic LLM safety (to our knowledge):} language/category-conditioned steering vectors extracted from contrastive safe/unsafe pairs.
    \item \textbf{SafeSteer-IN system:} a plug-in inference-time safety layer with no weight updates and low runtime overhead (targeting $<2$ms per-token steering overhead in microbenchmarks).
    \item \textbf{Preliminary evidence:} early steering results on Hindi/Tamil slices plus analyses of layer sensitivity and cross-linguistic vector alignment.
    \item \textbf{Reproducible artifacts:} open-source code, vector extraction/evaluation scripts, and planned release of dataset/vector assets under Apache 2.0 / CC BY 4.0.
\end{itemize}

\section{India Safety Taxonomy}
\subsection{Taxonomy Overview}
The intellectual center of SafeSteer-IN is an India-grounded taxonomy implemented in our dataset with 8 categories: communal and religious hate, caste discrimination, political misinformation, gender based violence, code mixed toxicity, anti minority sentiment, child safety, and financial scam. The taxonomy is paired with 6 language settings in our code file: Hindi (hi), Tamil (ta), Bengali (bn), Gujarati (gu), Marathi (mr), and Hinglish (hi-en).

\begin{table*}[t]
\small
\centering
\begin{tabularx}{\textwidth}{p{2.5cm} X X c}
\toprule
\textbf{Category} & \textbf{India-grounded description} & \textbf{Linguistic/cultural signal} & \textbf{In existing benchmarks?} \\
\midrule
communal and religious hate & Incitement targeting religious communities, riot justification, demonization in locally salient narratives. & Region-specific communal rhetoric and lexical cues not captured by generic hate templates. & Partial \\
caste discrimination & Harmful content around caste hierarchy, Dalit dehumanization, untouchability, anti-intercaste rights. & Caste slurs and status terms with no direct Western equivalents. & No \\
political misinformation & Election rumors, fabricated politician statements, policy misinformation in Indian contexts. & Local entities, election terms, and event-specific misinformation patterns. & Partial \\
gender based violence & Regionally grounded misogynistic harm, including honor-based violence narratives. & Vernacular coercion/threat forms and culturally specific justifications. & Partial \\
code mixed toxicity & Harmful content using script/language mixing to evade moderation. & Hinglish/Tanglish mixing, romanized abuse tokens, mixed orthography. & No \\
anti minority sentiment & Targeting marginalized communities in Indian settings (e.g., tribal, NE communities). & Community-specific slurs and stereotypes absent in global datasets. & Partial \\
child safety & Child marriage normalization, exploitation facilitation, legal-evasion framing. & Terms tied to local legal/social discourse and community practices. & Partial \\
financial scam & UPI/KYC/OTP scam scripting and social engineering in local languages. & India-specific fintech vocabulary (UPI, KYC, OTP) and scam templates. & No \\
\bottomrule
\end{tabularx}
\caption{SafeSteer-IN India safety taxonomy. The final column summarizes whether each harm type is directly represented in common multilingual/English safety benchmarks.}
\label{tab:taxonomy}
\end{table*}

This taxonomy is operationalized, not only conceptual: categories are used consistently in prompt templates, dataset construction, classifier labels, and steering slices.

\subsection{The Cultural Gap}
English-first moderation pipelines fail in India not only because of translation gaps, but because harm taxonomies themselves are incomplete for local social realities. Table~\ref{tab:gap} contrasts our categories against frequently used resources.

\begin{table}[t]
\small
\centering
\begin{tabular}{p{2.2cm}p{4.7cm}}
\toprule
\textbf{Benchmark} & \textbf{Missing vs. our 8 categories} \\
\midrule
ToxiGen \citep{hartvigsen-etal-2022-toxigen} & Direct coverage of caste discrimination, code-mixed toxicity, UPI/KYC scams: absent. \\
HatEval \citep{basile-etal-2019-semeval} & Focuses on targeted hate classes; limited coverage of India-specific caste, fintech fraud, and code-mixed abuse. \\
Indic safety resources (current) & Partial toxicity/hate coverage, but no unified taxonomy spanning all 8 culturally grounded categories with steering-ready contrastive pairs. \\
\bottomrule
\end{tabular}
\caption{Mini gap analysis motivating SafeSteer-IN for C3NLP: existing benchmarks under-cover culturally specific Indian harm patterns.}
\label{tab:gap}
\end{table}

\section{SafeSteer-IN System}
\subsection{Contrastive Activation Addition}
We build on CAA-style representation steering \citep{panickssery2024steering,zou2023representation}. For a layer $l$, a canonical steering direction is
\begin{equation}
\mathbf{v}_l = \mathbb{E}[\mathbf{h}^{l}_{\text{unsafe}}] - \mathbb{E}[\mathbf{h}^{l}_{\text{safe}}],
\end{equation}
where $\mathbf{h}^{l}$ denotes residual-stream activations (often at a selected token position). At inference, steering applies
\begin{equation}
\mathbf{h}^{l} \leftarrow \mathbf{h}^{l} - \alpha \mathbf{v}_l.
\end{equation}
In our implementation, this update is injected through a PyTorch forward hook, while vectors are normalized after extraction.

\subsection{Pipeline Architecture}
Figure~\ref{fig:pipeline} summarizes the full pipeline. First, an IndicBERT-based multi-task classifier predicts language and category labels from prompt text. It uses a shared IndicBERTv2 backbone and two heads (language and category), trained with joint cross-entropy loss in our classifer code. The runtime predictor outputs a risk score in $[0,1]$.

Second, the steering engine loads vectors from the dataset and applies slice-specific steering with configurable $\alpha$. The end-to-end orchestrator is the SafeSteer pipeline, which supports forced language/category overrides, thresholded steering, and fallback slice selection when an exact vector is unavailable.

\begin{figure*}[t]
\centering
\includegraphics[width=\linewidth]{figures/pipeline.pdf}
\caption{The SafeSteer-IN pipeline. An incoming prompt is first classified by the IndicBERT risk classifier (language + harm category + risk score). If risk score exceeds threshold, the steering engine loads the corresponding (language, category) vector and attaches a PyTorch forward hook that subtracts $\alpha \cdot \mathbf{v}$ from the residual stream at the target layer during generation. Optionally, Azure AI Content Safety provides an independent post-generation safety score.}
\label{fig:pipeline}
\end{figure*}

Per-slice design is central: one vector per (language, category, layer), with per-slice $\alpha$ calibration on held-out validation data.

\subsection{Dataset Construction}
Each contrastive record follows the schema used in our build dataset: 
\{prompt, safe\_response, unsafe\_response, language, category\}.
Seed pairs are provided in our data file across all categories/languages. For the current submission cycle, we use a synthetic-only expansion protocol: for each (language, category) slice, we expand seed pairs into larger contrastive sets with controlled LLM-based generation and template-backed fallback augmentation, then deduplicate and split into train/validation/test.

Importantly, augmentation is separated from the target steered model: synthetic/expanded data are generated externally and then used only to construct contrastive supervision for vector extraction. Current scale is approximately \PLACEHOLDER{N pairs per language-category slice} with \PLACEHOLDER{M total contrastive pairs} over \PLACEHOLDER{K steering slices}, and a held-out evaluation set of \PLACEHOLDER{E pairs}.

\section{Preliminary Experiments}
\subsection{Experimental Setup}
Primary experiments use Sarvam-1 (\PLACEHOLDER{parameter count}) with OpenHathi-7B as a secondary model. Candidate layers are selected from model config ranges (e.g., Sarvam-1 currently configured with layers 11--13 in \texttt{config.py}), and we save vectors per target layer for sensitivity analysis. For baseline-vs-steered evaluation, we construct a held-out harmful prompt benchmark with 20 prompts per (language, category) slice over 6 languages and 8 categories (total \PLACEHOLDER{960} prompts), generated with Claude Sonnet under fixed category/language constraints. The same prompt set is used for unsteered baselines and steered inference.

\subsection{Evaluation Metrics}
In code, aggregate metrics are implemented in our evaluation code through these safety metrics:
\begin{itemize}
    \item \textbf{SIR (Safety Intervention Rate):} fraction of harmful-prompt outputs that become safe after steering (reported as improvement over baseline safe counts).
    \item \textbf{FPR (False Positive Rate):} fraction of benign outputs negatively affected by steering.
    \item \textbf{Fluency:} perplexity-based preservation score (ratio of unsteered to steered perplexity on benign text; values near 1 indicate preservation).
    \item \textbf{Latency overhead:} mean additional generation latency attributable to steering hooks (reported in ms/token in final runs).
\end{itemize}
To improve reliability beyond rule-based heuristics, final safety labels are assigned by an LLM-as-a-judge protocol (Claude Sonnet 4.6): (i) we generate model outputs for each harmful prompt under unsteered and steered conditions, (ii) we score each output for harmfulness with a fixed judging rubric and structured output schema, and (iii) we compute slice-level harmful-output rates and safe-response rates from these judgments. We additionally run manual spot checks to validate judge consistency.

\subsection{Main Results}
\begin{table*}[t]
\small
\centering
\begin{tabular}{l l l c c c c}
\toprule
\textbf{Model} & \textbf{Language} & \textbf{Category} & \textbf{SIR} & \textbf{FPR} & \textbf{Fluency} & \textbf{Latency} \\
\midrule
Sarvam-1 & hi & communal\_religious\_hate & \PLACEHOLDER{value} & \PLACEHOLDER{value} & \PLACEHOLDER{value} & \PLACEHOLDER{ms/token} \\
Sarvam-1 & hi & caste\_discrimination & \PLACEHOLDER{value} & \PLACEHOLDER{value} & \PLACEHOLDER{value} & \PLACEHOLDER{ms/token} \\
Sarvam-1 & hi & code\_mixed\_toxicity & \PLACEHOLDER{value} & \PLACEHOLDER{value} & \PLACEHOLDER{value} & \PLACEHOLDER{ms/token} \\
Sarvam-1 & ta & communal\_religious\_hate & \PLACEHOLDER{value} & \PLACEHOLDER{value} & \PLACEHOLDER{value} & \PLACEHOLDER{ms/token} \\
Sarvam-1 & ta & caste\_discrimination & \PLACEHOLDER{value} & \PLACEHOLDER{value} & \PLACEHOLDER{value} & \PLACEHOLDER{ms/token} \\
Sarvam-1 & ta & code\_mixed\_toxicity & \PLACEHOLDER{value} & \PLACEHOLDER{value} & \PLACEHOLDER{value} & \PLACEHOLDER{ms/token} \\
OpenHathi-7B & hi & communal\_religious\_hate & \PLACEHOLDER{value} & \PLACEHOLDER{value} & \PLACEHOLDER{value} & \PLACEHOLDER{ms/token} \\
OpenHathi-7B & ta & caste\_discrimination & \PLACEHOLDER{value} & \PLACEHOLDER{value} & \PLACEHOLDER{value} & \PLACEHOLDER{ms/token} \\
\bottomrule
\end{tabular}
\caption{Preliminary slice-wise results. All numbers are placeholders pending final reruns. Target trend: high SIR with controlled FPR and small latency overhead.}
\label{tab:main-results}
\end{table*}

Preliminary trends suggest steering effectiveness is substantial on high-risk slices, with expected SIR above \PLACEHOLDER{70\%} and FPR below \PLACEHOLDER{15\%} in the strongest settings. We treat these as early indicators rather than finalized claims, since full six-language/eight-category runs are still ongoing.

\subsection{Layer Sensitivity Analysis}
\begin{figure}[t]
\centering
\includegraphics[width=\linewidth]{figures/layer_sensitivity.pdf}
\caption{SIR (\%) vs. transformer layer index for Sarvam-1 on Hindi communal hate and caste discrimination categories. Steering is most effective in the middle layers (approx. layers 10--14), consistent with prior findings on English models.}
\label{fig:layers}
\end{figure}

Figure~\ref{fig:layers} studies SIR as a function of steering layer. We observe a mid-layer peak around \PLACEHOLDER{layer interval}, while very early and very late layers underperform. This aligns with prior representation-engineering observations in English models, and provides initial evidence that similar mechanistic structure may transfer to Indic LLMs.

\subsection{Cross-Linguistic Steering Vector Alignment}
\begin{figure}[t]
\centering
\includegraphics[width=\linewidth]{figures/cosine_heatmap.pdf}
\caption{Pairwise cosine similarities between steering vectors for the communal\_hate category across evaluated languages (Hindi, Tamil) for Sarvam-1. Higher similarity suggests shared internal safety representations.}
\label{fig:cosine}
\end{figure}

We compute pairwise cosine similarities between vectors of the same category across languages (Figure~\ref{fig:cosine}). Our hypothesis is that if Indo-Aryan languages (Hindi, Bengali, Gujarati, Marathi) cluster more tightly than Tamil, then safety representations partly track language-family structure. Preliminary two-language runs are consistent with \PLACEHOLDER{observed trend}; full multilingual validation is pending. If confirmed, this has implications for cross-lingual transfer and parameter sharing in safety steering \citep{wang2025clas,soteria2025}.

\section{Limitations and Conclusion}
\section*{Limitations}
This work is explicitly preliminary. First, full coverage over 6 languages and 8 categories is ongoing; current quantitative slices are limited to a subset. Second, part of the contrastive data is synthetic or LLM-augmented, which may not fully match in-the-wild distributions. Third, formal large-scale native-speaker human evaluation has not yet been completed.

Fourth, alpha is currently calibrated per slice as a fixed scalar; adaptive alpha policies conditioned on risk score remain future work. Fifth, the IndicBERTv2 classifier used for routing is encoder-based and not instruction-tuned; it may rely on surface cues and can be brittle for adversarially phrased prompts. Relatedly, heuristic/automated evaluators can miss nuanced harms, motivating LLM-as-a-classifier and multi-judge evaluation designs in future iterations. Finally, model coverage remains limited, and cross-model generalization of vectors is not yet established.

\section*{Ethics Statement}
All harmful content in this study is used solely for safety research and red-teaming of mitigation mechanisms. No real user conversation logs were collected. Data construction is based on curated templates and public datasets; we plan controlled release with researcher-use agreements, usage restrictions, and explicit prohibition of misuse for harm generation. Because safety interventions can over-refuse or unevenly affect dialect communities, we treat deployment as assistive rather than fully autonomous and prioritize human oversight and culturally informed review.

\section*{Conclusion}
We introduced SafeSteer-IN, a culturally grounded inference-time safety steering system for Indic LLMs, centered on an India-specific harm taxonomy and CAA-based activation control. Our implementation demonstrates an end-to-end path from taxonomy and contrastive data construction to slice-specific steering and evaluation without model fine-tuning. Ongoing work extends evaluation across all 6 language settings and 8 categories with stronger human and multi-judge protocols.

\bibliography{refs}

\end{document}